\documentclass{amsart}
% For dealing with references we use the comment environment
\usepackage{verbatim}
\newenvironment{reference}{\comment}{\endcomment}
%\newenvironment{reference}{}{}
\newenvironment{slogan}{\comment}{\endcomment}
\newenvironment{history}{\comment}{\endcomment}

% For commutative diagrams we use Xy-pic
\usepackage[all]{xy}

% We use 2cell for 2-commutative diagrams.
\xyoption{2cell}
\UseAllTwocells

% We use multicol for the list of chapters between chapters
\usepackage{multicol}

% This is generally recommended for better output
\usepackage{lmodern}
\usepackage[T1]{fontenc}

% For cross-file-references

% Package for hypertext links:
\usepackage{hyperref}

% For any local file, say "hello.tex" you want to link to please

% Theorem environments.
%
\theoremstyle{plain}
\newtheorem{theorem}[subsection]{Theorem}
\newtheorem{proposition}[subsection]{Proposition}
\newtheorem{lemma}[subsection]{Lemma}

\theoremstyle{definition}
\newtheorem{definition}[subsection]{Definition}
\newtheorem{example}[subsection]{Example}
\newtheorem{exercise}[subsection]{Exercise}
\newtheorem{situation}[subsection]{Situation}

\theoremstyle{remark}
\newtheorem{remark}[subsection]{Remark}
\newtheorem{remarks}[subsection]{Remarks}

\numberwithin{equation}{subsection}

% Macros
%
\def\lim{\mathop{\mathrm{lim}}\nolimits}
\def\colim{\mathop{\mathrm{colim}}\nolimits}
\def\Spec{\mathop{\mathrm{Spec}}}
\def\Hom{\mathop{\mathrm{Hom}}\nolimits}
\def\Ext{\mathop{\mathrm{Ext}}\nolimits}
\def\SheafHom{\mathop{\mathcal{H}\!\mathit{om}}\nolimits}
\def\SheafExt{\mathop{\mathcal{E}\!\mathit{xt}}\nolimits}
\def\Sch{\mathit{Sch}}
\def\Mor{\mathop{\mathrm{Mor}}\nolimits}
\def\Ob{\mathop{\mathrm{Ob}}\nolimits}
\def\Sh{\mathop{\mathit{Sh}}\nolimits}
\def\NL{\mathop{N\!L}\nolimits}
\def\CH{\mathop{\mathrm{CH}}\nolimits}
\def\proetale{{pro\text{-}\acute{e}tale}}
\def\etale{{\acute{e}tale}}
\def\QCoh{\mathit{QCoh}}
\def\Ker{\mathop{\mathrm{Ker}}}
\def\Im{\mathop{\mathrm{Im}}}
\def\Coker{\mathop{\mathrm{Coker}}}
\def\Coim{\mathop{\mathrm{Coim}}}

% Boxtimes
%
\DeclareMathSymbol{\boxtimes}{\mathbin}{AMSa}{"02}

%
% Macros for moduli stacks/spaces
%
\def\QCohstack{\mathcal{QC}\!\mathit{oh}}
\def\Cohstack{\mathcal{C}\!\mathit{oh}}
\def\Spacesstack{\mathcal{S}\!\mathit{paces}}
\def\Quotfunctor{\mathrm{Quot}}
\def\Hilbfunctor{\mathrm{Hilb}}
\def\Curvesstack{\mathcal{C}\!\mathit{urves}}
\def\Polarizedstack{\mathcal{P}\!\mathit{olarized}}
\def\Complexesstack{\mathcal{C}\!\mathit{omplexes}}
% \Pic is the operator that assigns to X its picard group, usage \Pic(X)
% \Picardstack_{X/B} denotes the Picard stack of X over B
% \Picardfunctor_{X/B} denotes the Picard functor of X over B
\def\Pic{\mathop{\mathrm{Pic}}\nolimits}
\def\Picardstack{\mathcal{P}\!\mathit{ic}}
\def\Picardfunctor{\mathrm{Pic}}
\def\Deformationcategory{\mathcal{D}\!\mathit{ef}}

\setlength{\emergencystretch}{3em}
\begin{document}
\title[Stacks Project, Tag 0AXH]{428 --- Reduced hypersurfaces through prescribed divisors of a normal surface}
\maketitle
\noindent Copyright (C) 2005--2025 Johan de Jong. Stacks Project authors. Source: \href{https://stacks.math.columbia.edu/tag/0AXH}{Tag 0AXH}.

\noindent Editorial difficulty note (not author text). Difficulty H3: The author bounds the number of divisorial components under small perturbations and chooses a perturbation maximizing that number. Carefully selected symbolic-power valuations then force every component multiplicity to one without introducing new components, producing a reduced hypersurface through all prescribed divisors..

\noindent Editorial provenance note (not author text). History: excerpt from revision \texttt{a0444\allowbreak 6e57e\allowbreak c1fbc\allowbreak 252a8\allowbreak 71afc\allowbreak ec775\allowbreak 2fb28\allowbreak 07b14}, more-algebra.tex, lines 38137--38236. Reference presentation was adapted; original-author context and core support blocks were retained or added. The mathematical wording is unchanged. Licensed under GNU FDL 1.2 or later; no invariant sections or cover texts. See ../licenses/COPYING. Context blocks reproduce author definitions and supporting results; they are not additional counted problems.

\begin{lemma}
\label{lemma-radical-element}
Let $R$ be a Noetherian local normal domain of dimension $2$.
Let $\mathfrak p_1, \ldots, \mathfrak p_r$ be pairwise distinct
primes of height $1$. There exists a nonzero element
$f \in \mathfrak p_1 \cap \ldots \cap \mathfrak p_r$ such
that $R/fR$ is reduced.
\end{lemma}

\begin{proof}
Let $f \in \mathfrak p_1 \cap \ldots \cap \mathfrak p_r$ be a nonzero element.
We will modify $f$ slightly to obtain an element that generates a radical ideal.
The localization $R_\mathfrak p$ of $R$ at each height one prime
ideal $\mathfrak p$ is a discrete valuation ring, see
Algebra, Lemma \href{https://stacks.math.columbia.edu/tag/00PD}{lemma characterize dvr (Tag 00PD)} or
Algebra, Lemma \href{https://stacks.math.columbia.edu/tag/031S}{lemma criterion normal (Tag 031S)}.
We denote by $\text{ord}_\mathfrak p(f)$ the corresponding
valuation of $f$ in $R_{\mathfrak p}$. Let
$\mathfrak q_1, \ldots, \mathfrak q_s$
be the distinct height one prime ideals containing $f$.
Write $\text{ord}_{\mathfrak q_j}(f) = m_j \geq 1$ for each $j$.
Then we define $\text{div}(f) = \sum_{j = 1}^s m_j\mathfrak q_j$
as a formal linear combination of
height one primes with integer coefficients.
Note for later use that each of the primes $\mathfrak p_i$
occurs among the primes $\mathfrak q_j$.
The ring $R/fR$ is reduced if and only if
$m_j = 1$ for $j = 1, \ldots, s$. Namely, if $m_j$ is $1$ then
$(R/fR)\mathfrak q_j$ is reduced and
$R/fR \subset \prod (R/fR)_{\mathfrak q_j}$ as
$\mathfrak q_1, \ldots, \mathfrak q_j$ are the associated primes
of $R/fR$, see Algebra, Lemmas
\href{https://stacks.math.columbia.edu/tag/0311}{lemma zero at ass zero (Tag 0311)} and
\href{https://stacks.math.columbia.edu/tag/031T}{lemma normal domain intersection localizations height 1 (Tag 031T)}.

\medskip\noindent
Choose and fix $g$ and $N$ as in Lemma \href{https://stacks.math.columbia.edu/tag/0BWV}{lemma syspar (Tag 0BWV)}.
For a nonzero $y \in R$ denote $t(y)$ the number of primes minimal over $y$.
Since $R$ is a normal domain, these primes
are height one and correspond $1$-to-$1$ to the minimal primes of
$R/yR$ (Algebra, Lemmas \href{https://stacks.math.columbia.edu/tag/00KV}{lemma minimal over 1 (Tag 00KV)} and
\href{https://stacks.math.columbia.edu/tag/031T}{lemma normal domain intersection localizations height 1 (Tag 031T)}).
For example $t(f) = s$ is the number
of primes $\mathfrak q_j$ occurring in $\text{div}(f)$.
Let $h \in \mathfrak m^N$. By Lemma \href{https://stacks.math.columbia.edu/tag/0BWT}{lemma minprimespoly (Tag 0BWT)} we have
\begin{align*}
t(f + h) & \leq \text{length}_{R/(f + h)}(R/(f + h, g)) \\
& = \text{length}_R(R/(f + h, g)) \\
& = \text{length}_R(R/(f, g))
\end{align*}
see Algebra, Lemma \href{https://stacks.math.columbia.edu/tag/00IX}{lemma length independent (Tag 00IX)}
for the first equality.
Therefore we see that $t(f + h)$ is bounded independent of
$h \in \mathfrak m^N$.

\medskip\noindent
By the boundedness proved above we may pick
$h \in \mathfrak m^N \cap \mathfrak p_1 \cap \ldots \cap \mathfrak p_r$
such that $t(f + h)$ is maximal among such $h$. Set $f' = f + h$.
Given $h' \in \mathfrak m^N \cap \mathfrak p_1 \cap \ldots \cap \mathfrak p_r$
we see that the number $t(f' + h') \leq t(f + h)$.
Thus after replacing $f$ by $f'$ we may assume that for every
$h \in \mathfrak m^N \cap \mathfrak p_1 \cap \ldots \cap \mathfrak p_r$
we have $t(f + h) \leq s$.

\medskip\noindent
Next, assume that we can find an element $h \in \mathfrak m^N$ such that
for each $j$ we have $\text{ord}_{\mathfrak q_j}(h) \geq 1$ and
$\text{ord}_{\mathfrak q_j}(h) = 1 \Leftrightarrow m_j > 1$.
Observe that
$h \in \mathfrak m^N \cap \mathfrak p_1 \cap \ldots \cap \mathfrak p_r$.
Then $\text{ord}_{\mathfrak q_j}(f + h) = 1$
for every $j$ by elementary properties of valuations.
Thus
$$
\text{div}(f + h) = \sum\nolimits_{j = 1}^s \mathfrak q_j +
\sum\nolimits_{k = 1}^v e_k \mathfrak r_k
$$
for some pairwise distinct height one prime ideals
$\mathfrak r_1, \ldots, \mathfrak r_v$ and $e_k \geq 1$.
However, since $s = t(f) \geq t(f + h)$ we see that $v = 0$
and we have found the desired element.

\medskip\noindent
Now we will pick $h$ that satisfies the above criteria.
By prime avoidance (Algebra, Lemma \href{https://stacks.math.columbia.edu/tag/00DS}{lemma silly (Tag 00DS)})
for each $1 \leq j \leq s$ we can find an element $a_j \in \mathfrak q_j$
such that $a_j \not \in \mathfrak q_{j'}$ for $j' \not = j$
and $a_j \not \in \mathfrak q_j^{(2)}$. Here
$\mathfrak q_j^{(2)} = \{x \in R \mid \text{ord}_{\mathfrak q_j}(x) \geq 2\}$
is the second symbolic power of $\mathfrak q_j$.
Then we take
$$
h = \prod\nolimits_{m_j = 1} a_j^2 \times
\prod\nolimits_{m_j > 1} a_j
$$
Then $h$ clearly satisfies the conditions on valuations imposed above.
If $h \not \in \mathfrak m^N$, then we multiply by an element of
$\mathfrak m^N$ which is not contained in $\mathfrak q_j$ for all $j$.
\end{proof}
\end{document}
